\section{Round-sixteen positive closure: patched weak Harris theory and forced compressed memory}
\label{sec:r16-a4}

The complete resolved past is the state, but deterministic remote coordinates
are not minorized.  Contractivity is proved in a patched distance combining a
small-distance coupling with a Lyapunov penalty.  Sources form one fixed
Banach algebra.  The memory equation contains the unresolved initial forcing
unless the observable is resolved.

\subsection{The weighted history kernel}

Let \(\mathsf H=D(( -\infty,0],E)\) be the graph-completed past space and set
\[
 d_\beta(h,\tilde h)=\sum_{m\ge1}2^{-m}
  \bigl(1\wedge d_{J_1}(h|_{[-m,0]},\tilde h|_{[-m,0]})\bigr).
\]
Let \(W\ge1\) have compact sublevels and satisfy
\(PW\le aW+b\), \(a<1\), for the one-return past kernel \(P\).  Define the
small-distance cost
\[
 d_s(h,\tilde h)=\sqrt{d_\beta(h,\tilde h)
  \bigl(1+W(h)+W(\tilde h)\bigr)}
\]
and the patched metric
\[
 \widetilde d(h,\tilde h)=
 \min\{1,\delta^{-1}d_s(h,\tilde h)
       +\varepsilon|W(h)-W(\tilde h)|\}.
\]

\begin{theorem}[Patched weak-Harris contraction]
\label{thm:r16-a4-harris}
For suitable \(\delta,\varepsilon>0\), an integer \(n_0\), and
\(\theta<1\),
\[
 \mathcal W_{\widetilde d}(P^{n_0}(h,\cdot),
 P^{n_0}(\tilde h,\cdot))
 \le\theta\widetilde d(h,\tilde h).
\]
Consequently \(P\) has a unique invariant law on every communicating phase,
a spectral gap on the weighted Lipschitz space
\(\operatorname{Lip}_{\widetilde d,W}\), and exponential convergence in
\(\mathcal W_{\widetilde d}\).  No pure contraction is asserted in the raw
metric \(d_\beta\).
\end{theorem}

\begin{proof}
If \(d_\beta(h,\tilde h)\le\delta\) and both Lyapunov values are bounded, use
the common branch randomizer and maximal coupling of the next returned edge.
Summable past variation bounds the failure probability by
\(Cd_\beta(h,\tilde h)\).  On success the newest coordinates agree and every
old discrepancy is shifted one discounted level, giving a strict contraction
of \(d_s\).  On failure the distance is at most one, and the probability loss
is absorbed by the square-root factor in \(d_s\).

If one Lyapunov value is large, the drift inequality contracts the Lyapunov
part after a fixed number of steps.  On a common compact sublevel, a finite
connector set gives a \(d_s\)-small set: two trajectories enter a common
returned cylinder with uniformly positive probability.  Patching the local
coupling, the drift, and this small-set estimate gives the displayed
contraction.  The Wasserstein contraction theorem then gives uniqueness and
the weighted Lipschitz spectral gap.
\end{proof}

\subsection{One fixed analytic multiplier algebra}

Fix \(0<\eta<\eta_*\) and define
\[
 \mathscr A_\eta=
 \left\{V:\ 
 \|V\|_\eta:=\sup_h{|V(h)|\over1+\log W(h)}
 +\sup_{h\ne\tilde h}{|V(h)-V(\tilde h)|
  \over\widetilde d(h,\tilde h)(1+\log W(h)+\log W(\tilde h))}<\infty
 \right\}.
\]
Choose fixed weighted spaces
\(\mathcal X_+=\operatorname{Lip}_{\widetilde d,W^{1+2\eta}}\) and
\(\mathcal X_-=\operatorname{Lip}_{\widetilde d,W}\).

\begin{theorem}[Fixed Feynman--Kac chart]
\label{thm:r16-a4-fk}
There is \(r>0\) such that for \(\|V\|_\eta<r\),
\[
 P_Vf=P(e^Vf)
\]
is a holomorphic family in
\(\mathcal L(\mathcal X_+,\mathcal X_-)\), has a simple positive eigenvalue
\(e^{\lambda(V)}\), positive eigenfunction \(h_V\), and eigenmeasure
\(\nu_V\).  With \(\nu_V(h_V)=1\),
\[
 \widehat P_Vf=e^{-\lambda(V)}h_V^{-1}P_V(h_Vf)
\]
is a Markov kernel and the nonlinear operator
\[
 \mathcal E_VF=\log\widehat P_V(e^F)
\]
satisfies the exact tower law.
\end{theorem}

\begin{proof}
The fixed norm implies
\(e^V W\le C W^{1+\eta}\) and a weighted Lipschitz product estimate from
\(\mathcal X_+\) to \(\mathcal X_-\).  The exponential power series therefore
converges in one operator norm on the open ball.  The spectral gap from
Theorem~\ref{thm:r16-a4-harris} and analytic perturbation give the eigen-data.
Positivity and normalization make \(\widehat P_V1=1\).  The linear Markov
tower becomes the logarithmic nonlinear tower after exponentiation.
\end{proof}

\subsection{Poisson and enhanced martingale approximation}

For centered \(g\in\mathcal X_+\), define
\[
 h_g=\sum_{n\ge0}\widehat P_V^ng,
 \qquad
 D_{n+1}=h_g(H_{n+1})-
          \widehat P_Vh_g(H_n).
\]

\begin{theorem}[Uniform enhanced martingale approximation]
\label{thm:r16-a4-rough}
For every finite family \(g_1,\ldots,g_d\) in a bounded subset of
\(\mathcal X_+\), the rescaled additive functionals and their iterated
integrals converge, conditionally on every initial history with bounded
\(W\), to a Brownian rough path.  The covariance and area drift are
\[
 \Sigma=\mathbb E(D_1D_1^T),
 \qquad
 \Gamma={1\over2}
 \sum_{n\ge1}
 \mathbb E(g(H_0)\otimes g(H_n)-g(H_n)\otimes g(H_0)).
\]
The convergence is uniform on Lyapunov sublevels and survives suspension by
the A3 physical clock, including the pointed terminal residual.
\end{theorem}

\begin{proof}
The spectral gap makes the Poisson series converge in the weighted Lipschitz
norm.  The displayed \(D_{n+1}\) has zero conditional mean.  Its
\((2+\delta)\)-moment is uniformly bounded on Lyapunov sublevels, and the
conditional quadratic variation converges by the same spectral gap applied to
\(D D^T\).  Conditional Lindeberg follows by truncation and the Lyapunov
moment.  Burkholder estimates give tightness of first and second martingale
levels.  Expanding the coboundary remainder gives the absolutely convergent
antisymmetric Green--Kubo series \(\Gamma\).  The A3 clock theorem controls
the inverse clock and terminal residual uniformly, so the random time-change
mapping is continuous on the enhanced path space.
\end{proof}

\subsection{Unsmoothed suspension resolvent}

Let \(R_V(z)=(z-L_V)^{-1}\) be the continuous-time history resolvent.  The
entrance, complete-excursion, and terminal-residual operators are denoted by
\(A_V(z),K_V(z),B_V(z)\).  On the common multiplier chart they satisfy the
renewal identity
\[
 R_V(z)=R_V^{\rm res}(z)+A_V(z)(I-K_V(z))^{-1}B_V(z).
\]

\begin{theorem}[Raw renewal-resolvent theorem]
\label{thm:r16-a4-renewal}
On a strip \(\Re z>-\gamma\), the three renewal operators and their first two
source derivatives are meromorphic.  On every pole-free vertical contour,
\[
 \|A_V(z)\|+\|B_V(z)\|
 +\|(I-K_V(z))^{-1}\|
 \le C(1+|\Im z|)^{-2}.
\]
The bound is for the raw roof transform and follows from the
trace-compensated A2 theorem, not from multiplication by a smooth time test.
Consequently \(R_V(z)\) has integrable vertical bounds after its finitely many
resonant principal parts are removed.
\end{theorem}

\begin{proof}
Decompose a suspended path into entrance, complete returned excursions, and
one residual segment.  This gives the renewal identity first for \(\Re z\)
large.  The A2 raw Fourier estimate controls the complete-excursion transform
and its source derivatives on the whole vertical line.  Aperiodicity gives
invertibility away from isolated eigenvalue crossings.  On a pole-free
contour, the minimum singular value of \(I-K_V(z)\) is positive on the compact
low-frequency portion; at high frequency the raw norm is strictly below one.
Combining the two bounds gives a uniform inverse estimate.  Analytic Fredholm
continuation supplies the meromorphic strip and finitely many principal parts.
\end{proof}

\subsection{Common domains and forced compressed memory}

Let \(\mathcal H_V=L^2(\nu_V)\).  Transport \(\mathcal H_V\) to the reference
space \(\mathcal H_0\) by the unitary
\[
 U_Vf=f\sqrt{d\nu_V/d\nu_0}
\]
on the regular phase chart.  Let \(\mathcal R\subset D(L_0^2)\) be a fixed
finite-dimensional space and put
\(\mathcal R_V=U_V^{-1}\mathcal R\).  Thus
\(\mathcal R_V\subset D(L_V^2)\) and the orthogonal projections preserve the
common transported core.

Write \(P=P_{\mathcal R_V}\), \(Q=I-P\), and
\[
 C_V(z)=P(z-L_V)^{-1}P.
\]

\begin{theorem}[Forced compressed-resolvent identity]
\label{thm:r16-a4-memory}
For \(\Re z\) large, \(C_V(z)\) is strictly accretive and invertible.  Define
\[
 \widehat M_V(z)=zP-PL_VP-C_V(z)^{-1}.
\]
There is a unique causal operator-valued distribution \(M_V\) with this
Laplace transform.  For every \(A\in D(L_V)\),
\[
 x_A(t)=Pe^{tL_V}A
\]
satisfies
\[
 \dot x_A(t)=PL_VP x_A(t)
 +\int_0^tM_V(t-s)x_A(s)\,ds+\eta_A(t),
\]
where
\[
 \widehat\eta_A(z)
 =C_V(z)^{-1}P(z-L_V)^{-1}QA.
\]
The forcing vanishes exactly when \(A\in\mathcal R_V\).  The formula is
well-typed on the common operator core.
\end{theorem}

\begin{proof}
For \(y=(z-L_V)^{-1}x\) and \(x\in\mathcal R_V\), skew dissipativity of the
stationary Koopman generator gives
\(\Re\langle C_V(z)x,x\rangle=\Re z\|y\|^2>0\); finite-dimensionality gives
invertibility.  Taking Laplace transforms of \(Pe^{tL_V}A\), splitting
\(A=PA+QA\), and multiplying by \(C_V(z)^{-1}\) gives the displayed identity.
The large-\(z\) expansion on \(D(L_V^2)\) gives
\(\widehat M_V(z)=O(|z|^{-1})\), hence a unique causal distribution.  The
forcing formula is the remaining \(QA\) term.
\end{proof}

\begin{theorem}[Quantitative Schur descriptor theorem]
\label{thm:r16-a4-descriptor}
After adjoining to \(\mathcal R_V\) the finite residue ranges of all poles and
transmission zeros in \(\Re z\ge-\gamma_1\), there are \(c,C>0\) such that on
the shifted Bromwich contours
\[
 \sigma_{\min}(C_V(z))\ge c(1+|z|)^{-1},
 \qquad
 \|C_V(z)^{-1}\|\le C(1+|z|).
\]
The residual memory is an operator function satisfying
\[
 \|M_V(t)\|\le Ce^{-\gamma_1t},
 \qquad
 \int_0^\infty(1+t)\|M_V(t)\|dt<\infty.
\]
Instantaneous, resonant, transmission-zero, forcing, and residual components
are uniquely determined by the final Schur complement.
\end{theorem}

\begin{proof}
The raw renewal theorem gives a meromorphic finite-dimensional matrix with
polynomial vertical bounds and the expansion \(C_V(z)=z^{-1}P+O(z^{-2})\).
On a compact pole-free contour the smallest singular value has a positive
minimum.  At high frequency the expansion gives the stated lower bound.
Adjoining each finite residue and zero root space removes its principal factor
from the Schur complement; only finitely many occur in the closed strip.
Contour shifting and inverse Laplace transformation then give exponential
decay and the first moment.  Uniqueness follows from uniqueness of the partial
fraction decomposition and Laplace transform.
\end{proof}

\begin{theorem}[Round-sixteen A4 closure]
\label{thm:r16-a4-main}
The prepared history process has a patched weak-Harris spectral theorem, one
fixed analytic multiplier algebra, a centered enhanced martingale limit, an
unsmoothed suspension resolvent, and a fully forced, common-domain compressed
memory equation with quantitative inverse bounds.
\end{theorem}

\begin{proof}
Combine Theorems~\ref{thm:r16-a4-harris},
\ref{thm:r16-a4-fk}, \ref{thm:r16-a4-rough},
\ref{thm:r16-a4-renewal}, \ref{thm:r16-a4-memory}, and
\ref{thm:r16-a4-descriptor}.
\end{proof}
